\documentclass{amsart}

\usepackage{amsmath}
\usepackage{amssymb}
\usepackage{amsthm}
\PassOptionsToPackage{hyphens}{url}
\usepackage{hyperref}
\usepackage{booktabs}
\usepackage{array}
\usepackage{microtype}
\usepackage{seqsplit}
\emergencystretch=3em
\sloppy

\newtheorem{theorem}{Theorem}[section]
\newtheorem{proposition}[theorem]{Proposition}
\newtheorem{definition}[theorem]{Definition}
\newtheorem{corollary}[theorem]{Corollary}
\newtheorem{remark}[theorem]{Remark}

\title[An Auditable Scaffold for Beal's Conjecture]
{An Auditable Lean 4 Scaffold for Beal's Conjecture:\\
A Conditional Reduction and Its Level-26 Computational Evidence Base}

\author{David Fox}
\address{Opera Numerorum}
\email{david@zerobeacon.ai}

\subjclass[2020]{11D41, 11G05, 11F11, 68V20}
\keywords{Beal conjecture, Frey curve, $X_0(26)$, Lean 4, interactive theorem
proving, level lowering, axiom auditing, reproducibility}

\begin{document}

\begin{abstract}
We present two companion Lean~4 developments toward Beal's conjecture
($A^x+B^y=C^z$ with $\min(x,y,z)>2$ and $\gcd(A,B,C)=1$ implying a
contradiction) and describe the auditable scaffold architecture that
organizes both. The first development, \texttt{beal-conjecture}, reduces
the conjecture to exactly five explicit, ordinary theorem hypotheses via
a single assembly theorem; none of the five is a global \texttt{axiom}
declaration, and the reduction is machine-checked end to end. The second,
\texttt{beal-level-26-foundations}, is a companion computable evidence
base at level $26$: PARI/GP-verified two-descent, Selmer computations,
and cross-checks against LMFDB for the curves $26a1$ and $26b1$, all
reproducible from committed, hash-locked certificate data. We do
\emph{not} claim a proof of Beal's conjecture. Our contribution is
methodological as much as mathematical: a uniform Core/Wrapper/named-
hypothesis layering with automated axiom-footprint auditing
(\texttt{\#print axioms} on every declaration, enforced in continuous
integration), and an explicit translation table separating genuinely
axiom-free numerical bookkeeping from open mathematical obligations,
so that no internal label can be misread as more than it is. Every
claim in this note is anchored to a specific, versioned, DOI-archived
release; we point to those records rather than re-deriving their
contents inline.
\end{abstract}

\maketitle

\tableofcontents

\section{Introduction}

Beal's conjecture, advanced by Andrew Beal and first published as a prize
problem in 1997~\cite{beal1997}, asserts that for positive integers
$A,B,C,x,y,z$ with $x,y,z\ge 3$,
\[
A^x+B^y=C^z \implies \gcd(A,B,C)>1.
\]
No proof is known, and this note does not offer one. What we offer is an
\emph{auditable} pair of Lean~4 developments that make the logical
structure of the standard elliptic-curve approach to Beal completely
explicit, and a description of the audit methodology that keeps that
structure honest as the developments grow.

\subsection{Two repositories, one method}
The work is split across two repositories that play different roles.

\begin{itemize}
  \item \textbf{\texttt{beal-conjecture}} carries the \emph{conditional
  assembly}: a single theorem, \texttt{Beal.Final.ConditionalBealTheorem},
  that derives the unrestricted conjecture from five named mathematical
  premises, each an ordinary function argument rather than a global axiom.
  \item \textbf{\texttt{beal-level-26-foundations}} carries the
  \emph{computable evidence base}: explicit, reproducible, PARI/GP- and
  LMFDB-cross-checked computation at level $26$ (the curves $26a1$,
  $26b1$, their two-descent, and the associated Selmer data), together
  with a Galois-representation scaffold that narrows the shape of the
  remaining Tate/Ribet/modularity/Taylor--Wiles obligations without
  discharging them.
\end{itemize}

Both repositories are organized under the same audit discipline, which we
describe in Section~\ref{sec:audit} and regard as the primary
methodological contribution of this note. Detailed mathematical content
— discriminants, certificate hashes, exact Lean identifiers, and CI logs —
is not reproduced at length here; instead we cite the specific archived
Zenodo record and repository tag that carries each claim, so a reader can
verify the primary source directly (Section~\ref{sec:doi}).

\subsection{What is, and is not, claimed}
To avoid any ambiguity for a reader who does not open the Lean source: the
unrestricted statement of Beal's conjecture is not inhabited in either
repository. Section~\ref{sec:conditional} states exactly what
\texttt{ConditionalBealTheorem} establishes (a genuine but conditional
result). Section~\ref{sec:audit} states exactly what the level-26
repository's internal ``none''-labeled declarations establish, and
explicitly separates that from the conjecture itself, because the two are
easy to conflate from repository README prose alone.

\section{The Conjecture and the Conditional Reduction}
\label{sec:conditional}

\begin{definition}
A \emph{Beal solution} is a tuple $(A,B,C,x,y,z)\in\mathbf{Z}_{>0}^6$ with
$x,y,z>2$, $A^x+B^y=C^z$, and $\gcd(A,\gcd(B,C))=1$.
\end{definition}

\begin{definition}[\texttt{BealConjecture}, \texttt{lean/Beal/B01\_Def.lean}]
\label{def:bealconjecture}
No Beal solution exists.
\end{definition}

This is the literal, unrestricted statement: a universally quantified
proposition over all positive integers, with no numerical restriction. We
verified this directly against the current repository source rather than
assuming it from documentation.

The strategy for approaching Definition~\ref{def:bealconjecture} via
elliptic curves originates with Frey's observation~\cite{frey1986} that a
hypothetical primitive solution carries a geometric structure incompatible
with modularity. \texttt{beal-conjecture} v11.0.0 assembles this strategy
along a level-$26$ route through $X_0(26)$ and its Jacobian $J_0(26)$,
culminating in a single theorem.

\begin{theorem}[\texttt{Beal.Final.ConditionalBealTheorem}]
\label{thm:conditional}
Given
\begin{enumerate}
  \item \texttt{J0DecompositionSoundness\_26} — the soundness of the
  displayed decomposition of $J_0(26)$ used to locate the rational
  points of $X_0(26)$;
  \item \texttt{MwrankCertificateSoundness\_26} — the soundness of the
  supplied \texttt{mwrank} rank certificate for the relevant Mordell--Weil
  group;
  \item \texttt{FormalImmersionSoundness\_26} — the soundness of the
  formal-immersion criterion identifying the finite mod-$3$ matrix with
  the actual cotangent map at the boundary point;
  \item \texttt{FreyCurveExists} — existence of the Frey curve
  construction for a hypothetical primitive Beal solution; and
  \item \texttt{LevelLowering\_26} — the modularity and Ribet
  level-lowering data needed once that Frey curve is supplied,
\end{enumerate}
\texttt{BealConjecture} holds.
\end{theorem}

Each of the five premises in Theorem~\ref{thm:conditional} is stated as an
ordinary hypothesis of the theorem — not a \texttt{axiom} declaration —
so the dependency is visible in the type signature and cannot silently
propagate through unrelated declarations. The proof term constructs a
noncuspidal rational point on $X_0(26)$ from a hypothetical primitive
solution and contradicts the four-cusp characterization supplied by
premises~(1)–(3); we do not reproduce the twelve-line proof term here, as
it is short, self-contained, and directly readable in the cited release
(Table~\ref{tab:doi}).

\begin{remark}
Premise~(3) is the geometric compatibility gap: the finite mod-$3$ matrix
used in the formal-immersion criterion is derived from an explicit basis
change on normalized eigenform coefficients, but no theorem in the
repository identifies it with the literal Abel--Jacobi cotangent map.
Premise~(2)'s Selmer cardinality is proved only for an explicitly supplied
carrier group; its cohomological identification with the true Selmer
group of the curve remains external. These characterizations are taken
verbatim from the repository's own citation metadata, which we regard as
the authoritative, versioned statement of what remains open.
\end{remark}

\section{The Audit Architecture}
\label{sec:audit}

The methodological contribution of this note is the discipline that keeps
a large, multi-hundred-declaration Lean development honest about exactly
which classical axioms, and which named mathematical hypotheses, any given
result depends on. We describe the three-layer pattern used across both
repositories and the tooling that enforces it.

\subsection{Three layers}

\begin{itemize}
  \item \textbf{Core.} Files suffixed \texttt{\_Core.lean} carry zero
  imports from Mathlib and zero axioms beyond \texttt{propext}. Every
  Core declaration's type is the mathematical claim itself; every Core
  proof is verified with \texttt{\#print axioms} to return \texttt{[]}
  or \texttt{[propext]}.
  \item \textbf{Wrapper.} Ordinary (non-\texttt{\_Core}) files connect
  Core statements to Mathlib types (reals, \texttt{GCDMonoid}, field
  norms). The permitted axiom budget here is exactly
  $\{\texttt{propext}, \texttt{Classical.choice}, \texttt{Quot.sound}\}$
  — Mathlib's own classical footprint, never a project-introduced
  \texttt{sorry} or ad hoc axiom.
  \item \textbf{Named hypothesis.} A small number of declarations —
  Theorem~\ref{thm:conditional}'s five premises, and
  \texttt{beal-conjecture}'s historical \texttt{modularity\_hypothesis} —
  are declared exactly once as the boundary of current mathematical
  knowledge, and every downstream use is required (and CI-enforced) to
  route through that single named declaration rather than reintroducing
  the assumption informally elsewhere.
\end{itemize}

\subsection{Tooling}
Both repositories ship executable audits, run on every push:
\begin{itemize}
  \item a literal \texttt{grep -r "sorry"} over the Lean tree, required
  to return empty;
  \item \texttt{\#print axioms} on every Core declaration, checked against
  the permitted set above;
  \item a boundary check that a named hypothesis (e.g.\ the level-26
  repository's five-premise theorem, or \texttt{modularity\_hypothesis}
  in earlier drafts) is not accidentally depended upon by a declaration
  that claims to be axiom-free.
\end{itemize}

\subsection{A necessary disclaimer: ``none'' is a bookkeeping label, not a
proof of the conjecture}
\label{sec:disclaimer}

\texttt{beal-level-26-foundations} labels declarations \texttt{none} when
their \texttt{\#print axioms} output is empty (beyond the kernel), and
\texttt{propext only} when it is exactly $\{\texttt{propext}\}$. These
labels describe axiom footprint, and nothing else. In particular:

\begin{itemize}
  \item The repository's \texttt{BealForall} is a
  \texttt{structure}\ wrapping a single field,
  \texttt{frey\_beal\_forall\_none\_formula}, which itself unfolds to a
  conjunction of numeral identities and internal token equalities (for
  instance $2\cdot 13=26$, and the recorded discriminants of $26a1$ and
  $26b1$), each decidable by \texttt{rfl}. Declarations such as
  \texttt{BealForall\_real\_witness\_none} and
  \texttt{beal\_forall\_in\_kernel\_from\_beal\_forall\_none\_separated}
  are proofs of \emph{this} structure — a build-time consistency
  checklist — not of the quantified Beal statement.
  \item The quantified statement itself,
  \[
  \forall\, A\,B\,C\,m\,n\,p,\; 2<m\to 2<n\to 2<p\to A^m+B^n=C^p\to
  \gcd(A,\gcd(B,C))>1,
  \]
  named \texttt{beal\_forall\_from\_Is13Case\_sketch} in the repository,
  is proved \emph{uninhabited} in the ordinary sense of having no
  known term, and the repository carries an explicit lock,
  \texttt{original\_beal\_forall\_sketch\_type\_eq}, pinning down that
  this is the correct, unrestricted type.
\end{itemize}

We flag this explicitly because repository-level release notes describe
some of these declarations with language such as ``unconditional'' and
``BOTH none,'' which is accurate only in the narrow bookkeeping sense
above; read without this section, it is easy to mistake for a claim about
Definition~\ref{def:bealconjecture}. No such claim is made, here or in the
source.

\section{Computational Evidence Base at Level 26}
\label{sec:level26}

Independently of the bookkeeping discussed in Section~\ref{sec:disclaimer},
\texttt{beal-level-26-foundations} contains genuine, reproducible
computational number theory that narrows the shape of
Theorem~\ref{thm:conditional}'s premises at level $26$:

\begin{itemize}
  \item Two-descent on the elliptic curves $26a1$
  ($\Delta=-17576$) and $26b1$ ($\Delta=-1664$), computed via
  \texttt{cypari2}/PARI-GP, with the resulting Selmer, rank, and torsion
  data committed as SHA-256-hash-locked certificate JSON and cross-checked
  against LMFDB display data.
  \item An explicit finite scaffold recording the conductor arithmetic,
  cyclotomic character, and unramified/semistable/finite-flat conditions
  attached to the residual representation at the prime $13$, together
  with an explicit infinite family of auxiliary primes
  $q\equiv 1\pmod{13^n}$ used in the Taylor--Wiles patching argument's
  standard template.
\end{itemize}

We emphasize what this evidence does \emph{not} do: it does not discharge
any of the five premises of Theorem~\ref{thm:conditional}, and the
Galois-representation and Hecke-algebra structures referenced in the
scaffold are typed placeholders recording the \emph{shape} of the objects
Tate's algorithm, Ribet's theorem, and Taylor--Wiles patching would need
to supply — not verified instances of those theorems. We point to the
specific, versioned records (Table~\ref{tab:doi}) rather than reproducing
certificate contents inline; every certificate is reproducible from the
committed script and PARI-GP version recorded in that release.

\section{A Zero-Axiom Corollary: Beal Implies Fermat}

Independently of the conditional reduction, the following is a direct,
zero-axiom consequence of Definition~\ref{def:bealconjecture} at equal
exponents, verified against the current repository source.

\begin{corollary}[\texttt{beal\_implies\_fermat21\_core}]
\texttt{BealConjecture21Core} $\implies$ \texttt{FermatLastTheorem21Core}:
if no primitive Beal solution exists, then for all positive integers
$a,b,c$ and all $n>2$, $a^n+b^n\ne c^n$ whenever $\gcd(a,\gcd(b,c))=1$.
\end{corollary}

\begin{proof}
A Fermat instance at exponent $n>2$ with coprime $a,b,c$ is, by
definition, a Beal solution with $x=y=z=n$; the corollary is the
one-line term \texttt{fun h a b c n \_ hSol => h a b c n hSol}, checked
with \texttt{\#print axioms} to depend on nothing beyond the kernel.
\end{proof}

\section{Data and Reproducibility}
\label{sec:doi}

Table~\ref{tab:doi} anchors every claim above to a specific, archived,
versioned record. We follow the convention, used throughout both
repositories, of minting a Zenodo DOI on every tagged GitHub release and
recording it in that release's own documentation before any later release
references it.

\begin{table}[h]
\centering
\footnotesize
\caption{Key repositories, tags, and DOI records.}
\label{tab:doi}
\begin{tabular}{>{\raggedright\arraybackslash}p{3.0cm}>{\raggedright\arraybackslash}p{2.3cm}>{\raggedright\arraybackslash}p{6.6cm}}
\toprule
Repository / tag & DOI & Content \\
\midrule
\texttt{beal-conjecture} (concept) & \seqsplit{10.5281/zenodo.22041831} & Concept DOI; resolves to latest version. \\
\texttt{beal-conjecture} v11.0.0 & \seqsplit{10.5281/zenodo.22281075} & \texttt{ConditionalBealTheorem}, five named premises (Section~\ref{sec:conditional}). \\
\texttt{beal-level-26-\linebreak foundations} (concept) & \seqsplit{10.5281/zenodo.22272382} & Concept DOI; listed as \texttt{isSupplementTo} the conjecture repository above. \\
\texttt{beal-level-26-\linebreak foundations} v7.1.0 & \seqsplit{10.5281/zenodo.22632209} & Level-26 evidence base and audit-labeled bookkeeping (Sections~\ref{sec:audit}--\ref{sec:level26}). \\
\texttt{beal-level-26-\linebreak foundations} v7.1.1 & \seqsplit{10.5281/zenodo.22635221} & Documentation catch-up recording the v7.1.0 mint. \\
\bottomrule
\end{tabular}
\end{table}

Both repositories are MIT-licensed and publicly readable; every audit
script referenced in Section~\ref{sec:audit} is runnable from a fresh
clone with \texttt{lake build} followed by the repository's
\texttt{verify-scaffold.sh} (level-26 repository) or
\texttt{scripts/verify-all.sh} (conjecture repository).

\section{Remaining Obligations and Future Work}

The path from the present state to a genuine unconditional result runs
through discharging Theorem~\ref{thm:conditional}'s five premises in full
generality — most substantially, a Lean formalization of Tate's
algorithm~\cite{tate1975} sufficient to compute conductors for the general
Frey curve family, a formalization of Ribet's level-lowering
theorem~\cite{ribet1990} building on Mazur's irreducibility
criterion~\cite{mazur1978}, and the modularity input in the generality
Wiles and Taylor--Wiles established~\cite{wiles1995,taylorwiles1995}. None
of these exist in Mathlib at present; the Imperial College London FLT
project~\cite{flt2024} is formalizing closely related material for
Fermat's Last Theorem, and we regard that effort as the most direct
template for discharging the level-lowering premise here. The
level-26 evidence base's finite-flat, unramified, and Taylor--Wiles-family
scaffolding (Section~\ref{sec:level26}) is intended to narrow, incrementally,
the precise shape those formalizations would need to take at level $26$.

\section*{Acknowledgments}
The author thanks the Lean~4 and Mathlib communities for the
infrastructure that makes this level of audit possible, and the Imperial
College London FLT project for the formalization template and for making
the shape of the level-lowering gap explicit. The author maintains the
\textit{Opera Numerorum} formalization program and its accompanying
verification and reproducibility tooling, including the API
infrastructure at \href{https://zerobeacon.ai}{zerobeacon.ai}.

\begin{thebibliography}{99}

\bibitem{beal1997} D. Andrew Beal, \emph{The Beal conjecture and prize},
Notices Amer. Math. Soc. \textbf{44} (1997), no.~11, 1436--1437.

\bibitem{frey1986} Gerhard Frey, \emph{Links between stable elliptic
curves and certain Diophantine equations}, Ann. Univ. Sarav. Ser. Math.
\textbf{1} (1986), no.~1, 1--40.

\bibitem{glr2010} Enrique Gonz\'{a}lez-Jim\'{e}nez and Joan-Carles Lario,
\emph{Rational points on twists of $X_0(N)$}, Acta Arith. \textbf{143}
(2010), no.~2, 149--171.

\bibitem{flt2024} Kevin Buzzard et al., \emph{Fermat's Last Theorem in
Lean}, Imperial College London,
\url{https://github.com/ImperialCollegeLondon/FLT}, 2024.

\bibitem{mazur1978} Barry Mazur, \emph{Rational isogenies of prime
degree}, Invent. Math. \textbf{44} (1978), no.~2, 129--162.

\bibitem{ribet1990} Kenneth A. Ribet, \emph{On modular representations of
$\mathrm{Gal}(\overline{\mathbf{Q}}/\mathbf{Q})$ arising from modular
forms}, Invent. Math. \textbf{100} (1990), no.~2, 431--476.

\bibitem{tate1975} John Tate, \emph{Algorithm for finding the type of a
singular fiber in an elliptic pencil}, in: Modular Functions of One
Variable IV, Lecture Notes in Math. 476, Springer, Berlin, 1975, pp.~33--52.

\bibitem{taylorwiles1995} Richard Taylor and Andrew Wiles,
\emph{Ring-theoretic properties of certain Hecke algebras}, Ann. of Math.
(2) \textbf{141} (1995), no.~3, 553--572.

\bibitem{wiles1995} Andrew Wiles, \emph{Modular elliptic curves and
Fermat's last theorem}, Ann. of Math. (2) \textbf{141} (1995), no.~3,
443--551.

\bibitem{mathlib2020} The Mathlib Community, \emph{The Lean mathematical
library}, in: Proc. 9th ACM SIGPLAN Intl. Conf. Certified Programs and
Proofs (CPP 2020), ACM, New York, 2020, pp.~367--381.

\end{thebibliography}

\end{document}
